\chapter{Rappels de définitions}

\begin{itemize}
\it On considère sur une droite $(xy)$ trois points $O$, $A$, $B$ dans cet ordre et on construit le milieu $M$ de $[AB]$. Soit $OA=a$ et $OB=b$. \begin{enumerate}
\item Calculer à partir de $a$ et $b$ les longueurs 
$AM$, $MB$ et $OM$. 
\item Comment faut-il modifier les résultats si $O$ est entre $A$ et $M$ ? 
\end{enumerate}
\it Soit $B$ un point du segment $[AC]$. On désigne par $M$ le milieu de $[AB]$, par $N$ celui de $[BC]$ et on pose $AB = a$ et $BC = b$. \begin{enumerate}
\item Calculer à partir de $a$ et de $b$ la longueur $MN$. 
\item Soit $P$ le milieu de $[AC]$. Calculer $MP$ et $PN$. 
\end{enumerate}
\it Deux angles $\widehat{AOB}$ et $\widehat{BOC}$ sont adjacents et $[OM)$ est la bissectrice de l'angle 
$\widehat{BOC}$. \begin{enumerate}
\item Construire la figure sachant que $\widehat{AOB}=48^o$ et $\widehat{AOC}= 112^o$. Calculer la mesure de
l'angle $\widehat{AOM}$. 
\item Si $\alpha$ et $\beta$ désignent les mesures des angles $\widehat{AOB}$ et $\widehat{AOC}$, montrer que l'on a toujours $\widehat{AOM}= \frac12(\alpha+\beta)$.  
\end{enumerate}
\it Soient $[OM)$ et $[ON)$ les bissectrices de deux angles adjacents $\widehat{AOB}$ et $\widehat{AOC}$. 
\begin{enumerate}
\item Construire la figure pour $\widehat{AOB}= 64^o$ et $\widehat{AOC}= 108^o$.
\item Calculer la mesure de l'angle $\widehat{MON}$. Comparer le résultat à la mesure de $\widehat{BOC}$. 
\item Énoncer un théorème donnant la mesure de l'angle des bissectrices de deux angles adjacents de mesures $\alpha$ et $\beta$.
\end{enumerate}
\it On construit un angle $\widehat{AOB}= 72^o$, puis les angles droits $\widehat{AOA'}$ et $\widehat{BOB'}$ adjacents au premier. 
\begin{enumerate}
\item Calculer la mesure de l'angle $\widehat{A'OB'}$. Comment sont les deux angles $\widehat{AOB}$ et $\widehat{A'OB'}$ ? 
\item Montrer que les bissectrices $[OM)$ et $[ON)$ des angles $\widehat{AOB}$ et $\widehat{A'OB'}$ sont alignées.
\item Calculer l'angle des bissectrices des angles droits $\widehat{AOA'}$ et $\widehat{BOB'}$.
\end{enumerate}
\it Soit un angle $\widehat{AOB}$ de $54^o$. On construit les angles droits $\widehat{AOA'}$ et $\widehat{BOB'}$ non adjacents à l'angle $\widehat{AOB}$.\begin{enumerate}
\item Montrer que les angles $\widehat{AOB}$ et $\widehat{A'OB'}$ sont supplémentaires et qu'ils ont même bissectrice $[OM)$.
\item Démontrer que les angles $\widehat{AOB'}$ et $\widehat{A'OB}$ sont égaux. Calculer l'angle de leurs bissectrices.
\end{enumerate}
\it On construit trois angles successivement adjacents 
$\widehat{AOC} = 48^o$, $\widehat{COD}= 72^o$ et $\widehat{DOB}= 34^o$. Puis on trace les demi-droites $[OM)$, $[ON)$, $[OP)$ et $[OQ)$ bissectrices des angles $\widehat{AOC}$, $\widehat{AOD}$, $\widehat{BOC}$ et $\widehat{BOD}$. \begin{enumerate}
\item Calculer les angles $\widehat{MON}$ et $\widehat{POQ}$ et comparer leur valeur à celle de l'angle $\widehat{COD}$.
\item Montrer que les angles $\widehat{MOQ}$ et $\widehat{NOP}$ ont même bissectrice. 
\end{enumerate}
\it On considère sur un cercle les arcs consécutifs $\arc{AB} = 90^o$, $\arc{BC} = 68^o$ et $\arc{CD} = 90^o$. \begin{enumerate}
\item Calculer la mesure de l'arc $\arc{DA}$ qui termine le cercle. 
\item Construire au rapporteur le milieu $M$ de $\arc{BC}$ et le milieu $N$ de $\arc{DA}$. 
\item On plie la figure suivant la droite $(MN)$. Qu'en déduisez-vous pour le diamètre $[MN]$ et les cordes $[BC]$ et $[DA]$ ? 
\end{enumerate}
\it Soient $I$ et $J$ deux points d'une droite $(xy)$ et $A$ et $B$ deux points extérieurs situés d'un même côté de cette droite. Les cercles de centre $I$ et $J$ passant par $A$ se recoupent en $A'$, les cercles de centres $I$ et $J$ passant par $B$ se recoupent en $B'$. 
\begin{enumerate}
\item Montrer que si on plie la figure suivant $(xy)$, le point $A$ vient en $A'$ et le point $B$ en $B'$.
\item Montrer que $(xy)$ est médiatrice de $[AA']$ ainsi que de $[BB']$. 
\item La droite $(AB')$ coupe $(xy)$ en $M$. Montrer que $A'$, $M$ et $B$ sont alignés.
\end{enumerate}
\it \begin{enumerate}
\item Construire d'un même côté d'une droite $(xy)$ un polygone $F$ de six à huit côtés présentant un angle rentrant, puis le polygone $F'$ symétrique de $F$ par rapport à $(xy)$. 
\item Vérifier que les prolongements de deux côtés homologues $[BC]$ et $[B'C']$ par exemple se coupent sur $(xy)$ en un même point $I$. Comparer les angles $\widehat{BIx}$ et $\widehat{B'Ix}$.
\end{enumerate}
\it Soit un contour polygonal $ABCDEA'$ situé d'un même côté de la droite $(AA')$ et présentant un angle rentrant en $C$ ou en $D$. On désigne par $O$ le milieu de $[AA']$. \begin{enumerate}
\item Construire le contour $A'B'C'D'E'A$ symétrique du précédent par rapport au point $O$. Comparer les angles 
$\widehat{BC'D}$ et $\widehat{B'CD'}$. ainsi que les longueurs $BC'$ et $B'C$. 
\item Les segments $[BC']$ et $[CD']$ se coupent en $M$. Soit $M'$ le symétrique de $M$ par rapport à $O$. Démontrer que $M'$ se trouve à la fois sur $(CB')$ et sur $(DC')$. 
\end{enumerate}
\it On considère un angle droit $\widehat{AOB}$ et on dessine une ligne $F$, brisée ou comprenant des arcs de cercle, joignant $A$ à $B$, notée $AMNPQB$.\begin{enumerate}
\item Construire la figure $F_1$ et la figure $F_3$,
symétriques de $F$ par rapport à $(OA)$ et à $(OB)$.
Puis la figure $F_2$ symétrique de $F$ par rapport à $O$.
\item Montrer que $F_2$ et $F_1$ sont symétriques par rapport à la droite $(OB)$, tandis que $F_2$ et $F_3$ sont symétriques par rapport à $(OA)$, et que $F_1$ et $F_3$ sont symétriques par rapport à $O$. 
\item Quels sont les éléments de symétrie (centre et axes) de la figure totale formée par $F$, $F_1$, $F_2$ et $F_3$ ?
\end{enumerate}
\end{itemize}